% !TeX root = ../master.tex
%
% Frammento di ESEMPIO — riferimento di formato, non contenuto di un corso reale.
% Mostra come si scrive una lezione: apertura con \section, un ambiente
% tcbtheorem con la sua proof, e i due marker di revisione. Non va incluso in
% master.tex; serve solo da modello. Cancellalo pure quando non ti serve più.

\section{Titolo di esempio}

Testo introduttivo di esempio. Qui va il discorso della lezione, con la
notazione definita in \texttt{preamble.tex} (p.es. $\R$, $\N$) e in
\texttt{CLAUDE.md} del corso.

\begin{theorem}{Enunciato di esempio}{esempio-enunciato}
Sia $x \in \R$. Allora vale la proprietà di esempio $P(x)$.
\end{theorem}
\begin{proof}
Passaggio di esempio: da $P(x)$ segue $Q(x)$
\revisione{controllare questo passaggio: nell'immagine del prof il simbolo era poco leggibile}
e quindi la tesi.
\end{proof}

Nel punto in cui il professore aveva disegnato uno schema a mano si mette il
marker blu (che nelle lezioni vere è generato da \texttt{\textbackslash figura},
qui usato in forma grezza solo per mostrarne la resa):
\figuramancante{schema di esempio, da inserire}
